\section{Introduction}\label{sec:intro}

A human athlete is, from the point of view of an applied mathematician,
a non-linear stochastic dynamical system whose state is only ever observed
through a handful of noisy proxies. Heart rate, sleep duration, perceived
stress, and step counts are all loosely informative about the latent
fitness, fatigue, and autonomic condition of the body, but none of them
directly measures any of these quantities. The classical Banister model of
training response (Banister, 1991) treats the athlete as a deterministic
two-compartment system; the present notes work with a stochastic
three-compartment refinement, denoted FSA-v2, in which fitness, fatigue,
and an autonomic amplitude evolve as a coupled It\^o stochastic differential
equation driven by a single exogenous control --- the daily training
stimulus $\Phi(t)$.

Two questions follow naturally:

\begin{enumerate}
  \item \emph{Filtering.} Given a stream of multi-channel daily
        observations, what is the joint posterior distribution over the
        latent state $(B,F,A)$ and the static parameters $\theta$ of the
        SDE?
  \item \emph{Control.} Given the same posterior, what schedule
        $\Phi:[0,T]\to[0,\Phimax]$ maximises a chosen functional of the
        latent trajectory, subject to a soft constraint that fatigue
        $F(t)$ should not exceed a physiological ceiling $\Fmax$?
\end{enumerate}

Both questions sit at the intersection of Bayesian inference and stochastic
optimal control. They are usually treated as separate disciplines; one of
the contributions of the underlying code base is to express both within a
single sequential Monte Carlo framework. The filter is built from a
bootstrap particle filter, lifted to a sampler over static parameters by
the SMC${}^2$ algorithm (Chopin, Jacob and Papaspiliopoulos, 2013) and
operated in a rolling-window mode. The controller is built by recasting
the cost-minimisation problem as posterior inference under a tempered cost
likelihood, in the spirit of Toussaint (2009) and the more general
control-as-inference literature (Kappen, 2005; Levine, 2018). The two
machines are then composed in a model-predictive controller that re-plans
the next twelve hours of training stimulus from the most recent posterior.

These notes are organised so that the mathematical content can be read
linearly. Section~\ref{sec:fsa} defines the FSA-v2 SDE and observation
model in their G1-reparametrized form. Section~\ref{sec:fim} carries
out the Fisher-information-matrix analysis: an analytical derivation
of the three rank-deficient parameter pairs of the original
parametrization, the reparametrization that absorbs them, and the
empirical evidence that the resulting model is fully identifiable at
one-day windows. Section~\ref{sec:lyap} establishes stability of the
SDE: Jacobian-based local analysis of the deterministic skeleton, the
Stuart--Landau bifurcation locus, and a quadratic Lyapunov function
satisfying a Foster--Lyapunov drift inequality, from which exponential
ergodicity follows. Section~\ref{sec:num-stab} bridges from continuous-time
theory to numerical practice: it asks whether the fifteen-minute outer
step of the explicit Euler--Maruyama solver actually used in the code
base is small enough for mean-square stability, finds the answer to be
``yes, with three orders of magnitude of safety margin'', and identifies
the resolution of the sub-daily $\Phi$-burst as the binding constraint
should one wish to enlarge the step to speed up long-horizon experiments.
Section~\ref{sec:pf} develops the bootstrap
particle filter. The SMC${}^2$ outer loop and its rolling-window
variant are presented in Section~\ref{sec:smc2}.
Section~\ref{sec:control} formulates the control problem as a
posterior-inference problem. Section~\ref{sec:mpc} is the extended
core of the document: it develops the receding-horizon principle,
describes the FSA-v2 model-predictive controller in both mathematical
and implementation terms, and sketches a linear--quadratic--Gaussian
approximation that could be added as a cheap baseline.
Section~\ref{sec:results} reports the two headline empirical results
of the code base: an open-loop gain of $+28\%$ in mean autonomic
amplitude over a constant-stimulus baseline on an $84$-day horizon,
and a closed-loop gain of $+17\%$ in a fourteen-day rolling
controller. Section~\ref{sec:discussion} closes by reviewing the
identifiability and stability analyses just completed, and by
identifying the LQG / Riccati alternative as the principal piece of
remaining theoretical work.
